% ===================================================================
% Section 1: Introduction and motivation
% ===================================================================
\section{Introduction and motivation}

Markowitz mean--variance selection \cite{markowitz1952} is the foundation of
quantitative portfolio construction, but in its pure form it is fragile. The
global minimum-variance (GMV) portfolio,
\[
  \w_{\GMV}=\frac{\Sig^{-1}\one}{\one^\top\Sig^{-1}\one},
\]
depends on the \emph{inverse} of the covariance matrix, and in practice $\Sig$ is
replaced by a sample estimate $\Sighat$ formed from a finite window of returns.
Estimation error in $\Sighat$ is amplified by the inversion: the smallest
eigenvalues of $\Sighat$---the least well-estimated directions---dominate
$\Sighat^{-1}$, producing weights that are extreme, concentrated in a few
assets, unstable across rebalances, and expensive to trade
\cite{brittenjones1999}. When the number of assets $N$ approaches the number of
observations $T$, $\Sighat$ becomes ill-conditioned or singular and the problem
is acute.

The standard responses regularize either the estimate or the weights. Ridge
regularization (diagonal loading) replaces $\Sighat$ with $\Sighat+\lambda I$;
Ledoit--Wolf shrinkage pulls $\Sighat$ toward a structured target
\cite{ledoitwolf2004}; norm constraints on $\w$ stabilize the solution
\cite{demiguel2009norm,brodie2009}; risk-budgeting methods equalize risk
contributions \cite{maillard2010}; and the naive $1/N$ portfolio, which uses
\emph{no} estimate at all, is a famously hard benchmark to beat out of sample
\cite{demiguel2009}. Each of these is, in effect, a statement about the
\emph{spectrum} of the risk model: do not let the optimizer exploit near-flat
directions that are really just estimation noise.

\emph{Optimal experimental design} is the branch of statistics that formalizes
exactly this concern. A design chooses how to allocate observational effort so
that the resulting parameter estimates are as precise as possible, and it does so
by scalarizing an information matrix through the classical A-, D-, and
E-optimality criteria \cite{bergerwong,pukelsheim2006,kiefer1960}. The central
observation of this paper is that portfolio weights on the budget simplex play
the same role as a design measure, and that the A/D/E criteria supply a
principled, interpretable, and \emph{convex} family of regularizers---each
targeting a different, named failure mode of the raw GMV portfolio.

\subsection{Contribution}
Our contributions are as follows.

\begin{enumerate}[leftmargin=1.6em,itemsep=0.3em]
  \item \textbf{A unified convex formulation.} We formulate a single optimizer
    (Section~\ref{sec:criteria}) that combines mean--variance risk, ridge
    regularization, and A/D/E design penalties over a composable constraint set
    of budget, box, and group restrictions.

  \item \textbf{A regression interpretation of the design points.} We give the
    design vectors $\vv_i$ a genuine factor-regression meaning
    (Section~\ref{sec:vi}): they are the factor loadings of asset $i$, the target
    parameter is the vector of factor premia, and the weight-dependent
    information matrix $\M=B^\top\diag(\w)B$ measures the precision with which the
    \emph{portfolio} identifies the common factors.

  \item \textbf{A theory of the choice of design matrix.} Section~\ref{sec:Vchoice}
    analyses how the criteria depend on $V$. We show that the D-criterion is
    invariant under any invertible reparametrization of the factor space, whereas
    the A- and E-criteria are invariant only under orthogonal changes of basis
    (Theorem~\ref{thm:invariance}), and that under rescaling of $V$ the D-criterion
    is invariant in the argmin while A and E are not
    (Proposition~\ref{prop:scaling}). This has a concrete modelling consequence
    for statistically estimated factors, which are identified only up to rotation,
    and it explains why a normalization of $V$ is \emph{required} for the penalty
    strengths to be comparable across universes.

  \item \textbf{The rank-deficient case.} The familiar reductions of the D- and
    A-criteria to a log-barrier and a weighted inverse barrier hold only when $V$
    is square and nonsingular. For a genuine factor model with $K<N$ we give the
    correct Cauchy--Binet representation of the D-criterion
    (Theorem~\ref{thm:cauchybinet}) and a necessary and sufficient support
    condition for the information matrix to be nonsingular
    (Corollary~\ref{cor:spanning}), recovering the classical barrier forms as the
    special case $K=N$ (Corollary~\ref{cor:KN}).

  \item \textbf{Structural theory.} Section~\ref{sec:theory} establishes standing
    assumptions, convexity with full proofs, existence and uniqueness, exact
    recovery of the GMV and ridge-GMV portfolios, and limiting solutions as the
    penalties grow without bound.

  \item \textbf{Monotone comparative statics and stability.} We prove that each
    criterion evaluated at the optimum is monotone in its own penalty strength
    (Theorem~\ref{thm:monotone}), which explains the monotone sweeps observed in
    our experiments; and we prove a stability theorem
    (Theorem~\ref{thm:stability}) showing that the solution map
    $\Sighat\mapsto\w^\star$ is Lipschitz with constant inversely proportional to
    the strong-convexity modulus, which the penalties increase. This turns the
    paper's headline empirical finding---an order-of-magnitude turnover
    reduction---into a provable consequence of the formulation.

  \item \textbf{Empirical evaluation.} We evaluate on synthetic data, in a
    high-dimensional simulation with $N/T\to1$, and on a real universe of US
    sector ETFs (Sections~\ref{sec:computation}--\ref{sec:applications}),
    reporting both the gains and the honest trade-offs.
\end{enumerate}

\subsection{Organization}
Section~\ref{sec:setup} sets up the portfolio problem, the design bridge, the
factor interpretation of the design points, and the theory of the choice of $V$.
Section~\ref{sec:criteria} defines the three regularized portfolios and their
economic meaning. Section~\ref{sec:theory} develops the structural theory.
Section~\ref{sec:computation} covers computation, tuning, and the simulation
studies; Section~\ref{sec:applications} the real-data application; and
Section~\ref{sec:discussion} discusses when each criterion is useful. Longer
proofs are collected in the appendix.
